\chapter{Example Programs}
\label{ch:example-programs}

This appendix contains sample programs that demonstrate several features of the CPU.

\section{Byte Sieve}

This program computes primes below 1000 using the Sieve of Eratosthenes, storing one byte
per candidate number in a memory-mapped array rather than in a register or the stack. It is
a compact illustration of addressing-register setup (\reg{ah}/\reg{al} and \reg{sh}/\reg{sl})
and of using \mnemonic{STOR} and \mnemonic{LOAD} with an indexed effective address to walk an
array. The two nested loops starting at labels \texttt{:2} and \texttt{:3} (lines 35--50) are
the core of the sieve: the outer loop advances the candidate, and the inner loop marks its
multiples as composite.

% These example programs are too long to fit a non-breaking [H] float, so they use
% \lstinputlisting's own (page-breaking) caption directly rather than the "example" float
% used elsewhere in the manual, while still stepping the shared example counter and adding a
% matching "List of Examples" entry so they appear there too, consistently numbered.
\refstepcounter{example}
\label{ex:byte-sieve}
\addcontentsline{loe}{example}{\protect\numberline{\theexample}Byte sieve}
\lstinputlisting[caption={Example~\theexample: Byte sieve}]{../../examples/byte-sieve/byte-sieve.sasm}

\section{Faults}

This pair of programs exercises every fault type the CPU can raise and the vector table that
routes them to handlers. \texttt{faults.sasm} deliberately triggers, in sequence, a bus fault,
an alignment fault, two segment-overflow faults (one from an effective-address calculation,
one from program-counter overflow), an instruction-trace fault, a bus-protection fault, an
invalid-opcode fault (which itself triggers a nested double fault), a privilege violation, and
a level-five interrupt. Each handler, defined from line 144 onward, increments a distinct
register so a reader can confirm which faults fired and in what order. The segment-overflow
handler (lines 167--204) is worth particular attention: it reads exception metadata with
\mnemonic{ETFR} to distinguish an instruction-fetch overflow from an effective-address
overflow before deciding how to fix up the return address.

\refstepcounter{example}
\label{ex:faults}
\addcontentsline{loe}{example}{\protect\numberline{\theexample}Fault handling}
\lstinputlisting[caption={Example~\theexample: Fault handling}]{../../examples/faults/faults.sasm}

\texttt{faults-high.sasm} is a second, separately assembled segment that the program-counter
overflow test in Example~\ref{ex:faults} jumps into; it exists only to let that fault trigger
without the CPU executing the vector table sitting at address 0x0000 in the main segment. It
jumps back into \texttt{faults.sasm} through the trampoline at 0xFFF0.

\refstepcounter{example}
\label{ex:faults-high}
\addcontentsline{loe}{example}{\protect\numberline{\theexample}High-address fault handling}
\lstinputlisting[caption={Example~\theexample: High-address fault handling}]{../../examples/faults/faults-high.sasm}

\section{Hardware Exception}

This example shows a device-driven interrupt handled through the serial device: the program
prints a help message, waits on \mnemonic{WAIT} for a hardware exception, and exits once the
handler reports that the expected character has arrived in \reg{r7}. The exception handler
itself is installed at vector 0x0080 and defined from line 137 onward; it saves the link
register, calls \texttt{read\_pending\_byte} and \texttt{write\_pending\_byte} in turn, and
restores the link register before returning with \mnemonic{RETE}. Those two subroutines show
the poll-and-acknowledge pattern used to drain and refill the serial device's receive and
send buffers one byte at a time.

\refstepcounter{example}
\label{ex:hardware-exception}
\addcontentsline{loe}{example}{\protect\numberline{\theexample}Hardware exception}
\lstinputlisting[caption={Example~\theexample: Hardware exception}]{../../examples/hardware-exception/hardware-exception.sasm}

The data segment below supplies the two zero-terminated, word-per-character messages
(\texttt{help\_message} and \texttt{exit\_message}) that Example~\ref{ex:hardware-exception}
prints at startup and on exit.

\refstepcounter{example}
\label{ex:hardware-exception-data}
\addcontentsline{loe}{example}{\protect\numberline{\theexample}Hardware exception data}
\lstinputlisting[caption={Example~\theexample: Hardware exception data}]{../../examples/hardware-exception/data.sasm}

\section{Software Exception}

This short program demonstrates a software-raised exception with \mnemonic{EXCP} and confirms
that nesting is blocked while a handler is running. It sets \reg{r1} to 1, raises exception
0x40, and expects \reg{r1} to hold 0xFF once the handler at line 26 has run. That handler
itself raises a second exception, 0x41, from inside the handler; because the handler at vector
0x0500 (line 37) sets \reg{r1} to 0xBB only if it runs, the final value of \reg{r1} shows
whether the nested exception was correctly ignored.

\refstepcounter{example}
\label{ex:software-exception}
\addcontentsline{loe}{example}{\protect\numberline{\theexample}Software exception}
\lstinputlisting[caption={Example~\theexample: Software exception}]{../../examples/software-exception/software-exception.sasm}

\section{Store Load}

This program is a self-checking test of indexed and immediate-offset memory access, followed
by an exhaustive exercise of every \mnemonic{LJMP}/\mnemonic{LJSR} addressing form. The store
and load block (lines 8--19) round-trips values through memory using an immediate offset and a
register offset from an address-register pair. The remainder of the listing chains together
register-offset, immediate-offset, and plain register forms of \mnemonic{LJMP} and
\mnemonic{LJSR}; every path not on the intended route lands on an \texttt{LJMP s}/\texttt{LJSR
s} instruction commented \texttt{FAIL!}, so a working implementation reaches the \texttt{:pass}
label at the end while a faulty one jumps into the "forbidden zone" and fails at \texttt{:fail}.

\refstepcounter{example}
\label{ex:store-load}
\addcontentsline{loe}{example}{\protect\numberline{\theexample}Store/load}
\lstinputlisting[caption={Example~\theexample: Store/load}]{../../examples/store-load/store-load.sasm}
